\begin{problem}[Warmup]{(a)}
    Compute the following products (you should end up with a disjoint cycle representation): \[(134)(2457), \quad (12)(13)(14), \quad {(1567)}^3{(213)}^{-1}.\]
\end{problem}
\unsolved

\begin{problem}[Warmup]{(b)}
    Factor the following permutation as a product of disjoint cycles: \[\begin{pmatrix} 1 & 2 & 3 & 4 & 5 & 6 & 7 & 8 & 9 \\ 5 & 9 & 1 & 4 & 3 & 7 & 8 & 6 & 2 \end{pmatrix}.\] What is the order of this permutation?
\end{problem}
\unsolved

\begin{problem}[Warmup]{(c)}
    Give a permutation with three non-trivial orbits, one with 3 elements and two with 4 elements each. What is the order of the permutation you have constructed?
\end{problem}
\unsolved

\begin{problem}[Warmup]{(d)}
    Two non-identity permutations with disjoint supports necessarily commute, but the converse of this statement is false; give an example. \emph{Hint:} \emph{All} groups have non-trivial abelian subgroups!
\end{problem}
\unsolved

\begin{problem}[Warmup]{(e)}
    Does the set of 2-cycles in $S_4$ (together with the identity) form a subgroup? What about 3-cycles, or 4-cycles?
\end{problem}
\unsolved

\begin{problem}[Warmup]{(f)}
    In what ways are the Factorization Theorem for Permutations 30.3 and the Fundamental Theorem of Arithmetic 2.8 and 2.9 similar? In what ways are they different? (For the latter, consider whether cycles are \emph{irreducible}.)
\end{problem}
\unsolved

\begin{problem}{1}
    Compute explicitly the cyclic subgroups of $S_7$ generated by the following
    permutations:
    \[(357), \quad (14)(256), \quad (123)(456).\]
\end{problem}
\unsolved

\begin{problem}{2}
    Prove the uniqueness part of Factorization Theorem 30.3: The factorization of a
    permutation into disjoint cycles is unique, up to order.
\end{problem}
\unsolved

\begin{problem}{3}
    If $\alpha \in S_n$ is a cycle, is $\alpha^2$ a cycle also? Give an example when this
    is true and an example when this is false. Now characterize cycles where this is true
    and cycles where this is false.
\end{problem}
\unsolved

\begin{problem}{4}
    Give the order of each of the following permutations: $(123)(456)$, $(123)(4567)$,
    $(123)(45)$.
\end{problem}
\unsolved

\begin{problem}{5}
    Determine all possible orders of a product of two 3-cycles.
\end{problem}
\unsolved

\begin{problem}{6}
    Repeat Exercise 5 for the product of two 4-cycles.
\end{problem}
\unsolved

\begin{problem}{7}
    Suppose the permutation $\alpha = \beta_1\beta_2\cdots\beta_k$, the product of
    disjoint cycles where the order of cycle $\beta_i$ is $m_i$. Show that the order of
    $\alpha$ is the least common multiple of the $m_i$.
\end{problem}
\unsolved

\begin{problem}{8}
    \begin{enumerate}
        \item[(a)] There are four distinct disjoint cycle structures for non-identity
        elements of $S_4$. Name them.
        \item[(b)] The group of symmetries of the cube is isomorphic to $S_4$, since it can
        be viewed as the group of permutations of the diagonals of the cube. Describe what
        the four disjoint cycle structures of part a mean geometrically.
    \end{enumerate}
\end{problem}
\unsolved

\begin{problem}{9}
    Let $n$ be a positive integer, and $k, j$ fixed integers with $1 \le k, j \le n$. Let
    $G_k$ and $G_j$ be the stabilizer subgroups of $S_n$ considered in Exercise 29.6.
    Given $\alpha \in G_k$, define
    \[\Phi : G_k \to G_j\]
    by setting $\Phi(\alpha) = (kj)\alpha(kj)$ (where here $(kj)$ is a 2-cycle). Prove that
    $\Phi$ is a group isomorphism. (Note that an important verification is to check that
    $\Phi(\alpha) \in G_j$.)
\end{problem}
\unsolved

\begin{problem}{10}
    We generalize Exercise 9. Let $K$ and $J$ be subsets of $N = \{1, 2, 3, \ldots, n\}$,
    with the same number of elements. Let $G_K$ and $G_J$ be the stabilizer subgroups of
    $S_n$, discussed in Exercise 29.7. Since $K$ and $J$ have the same number of elements,
    there exists a one-to-one and onto function $\beta : K \to J$, and similarly there
    exists a one-to-one and onto function $\gamma : N \setminus K \to N \setminus J$.
    \begin{enumerate}
        \item[(a)] Define a function $\mu$ by setting $\mu(m) = \beta(m)$ if $m \in K$, and
        $\mu(m) = \gamma(m)$ if $m \in N \setminus K$. Argue that $\mu \in S_n$.
        \item[(b)] Given $\alpha \in G_K$, define $\Phi : G_K \to G_J$ by setting
        $\Phi(\alpha) = \mu\alpha\mu^{-1}$. Prove that this is a group isomorphism.
    \end{enumerate}
\end{problem}
\unsolved
